First the algebra of conjugation by $\sigma_r$. For $B$ in the plane pair: $\sigma_r\,\sigma_r\,\sigma_r = \sigma_r$ since $\sigma_r^2 = 1$; and $\sigma_r (I\sigma_r)\sigma_r = I\,\sigma_r^3 = I\sigma_r$ since $I$ commutes with even elements. For each of the four mixed planes, $B$ anticommutes with $\sigma_r$: writing $B = \gamma_a\gamma_b$ with exactly one index in $\{0,1\}$, the product $\sigma_r B$ reorders by two transpositions against one shared and one distinct generator, producing a net sign $-1$, so $\sigma_r B = -B\sigma_r$ and $\sigma_r B \sigma_r = -B\sigma_r^2 = -B$. Hence $-\tfrac12(B + 3\sigma_r B\sigma_r)$ equals $-2B$ on the plane pair and $+B$ on the mixed planes, which is $1 - 3\Pi$. Fact \claimref{interior.imported-schwarzschild-curvature-spectrum} says the curvature is diagonal on exactly these planes with eigenvalues $2M/r^3$ (doubled magnitude, on the pair) and $-M/r^3$ (mixed); since $1 - 3\Pi$ has eigenvalues $-2$ on the pair and $+1$ on the mixed planes, the operator with those eigenvalues is exactly $-\tfrac{M}{r^3}(1 - 3\Pi)$: so $R = -(M/r^3)S$ in the conventions of Fact \claimref{interior.imported-schwarzschild-curvature-spectrum}, fixed once in Section \ref{sec:setting}; only the eigenvalue ratio $2 : (-1)$ and the identity of the distinguished plane pair enter any later argument, and both are convention-free. The identity: $S^2 = (1 - 3\Pi)^2 = 1 - 6\Pi + 9\Pi^2 = 1 + 3\Pi$ since $\Pi^2 = \Pi$, so $S^2 + S = 2$. Spectrum and trace: eigenvalue $1 - 3 = -2$ on the two-dimensional range of $\Pi$ and $+1$ on its four-dimensional kernel; trace $= 4 - 4 = 0$.
